\ifdefined\gkver\else\def\gkver{student}\fi
\documentclass[\gkver]{gaokaozhenti}
\begin{document}

\chapter{2026年山东}

\begin{enumerate}
\item
%% number: 1
%% paperName: 2026年山东高考第 1 题 3 分
%% typeId: 1
%% type: 单选题
%% chapter: 原子和原子核
%% point: 人工核反应
%% method: 无
%% score: 3
%% degree: 950
%% duplicateId: 0
%% body:
我国钍资源丰富并成功实现了钍-铀核燃料转换，开辟了核燃料供应的新途径。转换过程的中间核素 $\ce{^{233}_{91}Pa}$ 可能会经历两种核反应，反应式为 $\ce{^{233}_{91}Pa + X -> ^{234}_{91}Pa}$，$\ce{^{233}_{91}Pa -> ^{233}_{92}U + Y}$，则 X、Y 分别是 \xzanswer{C}
\fourchoices[answer=C]
{$\ce{^{0}_{-1}e}$、$\ce{^{1}_{0}n}$}
{$\ce{^{1}_{0}n}$、$\ce{^{1}_{1}H}$}
{$\ce{^{1}_{0}n}$、$\ce{^{0}_{-1}e}$}
{$\ce{^{0}_{-1}e}$、$\ce{^{1}_{1}H}$}
%% answer: C
%% memo:
\memoanswer{
根据核反应的质量数和电荷数守恒可知，X 的质量数为 $1$，电荷数为 $0$，为中子 $\ce{^{1}_{0}n}$；Y 的质量数为 $0$，电荷数为 $-1$，则电子 $\ce{^{0}_{-1}e}$。

故选 C。
}

\item
%% number: 2
%% paperName: 2026年山东高考第 2 题 3 分
%% typeId: 1
%% type: 单选题
%% chapter: 直线运动
%% point: 运动图象
%% method: 无
%% score: 3
%% degree: 750
%% duplicateId: 0
%% body:
如图~\ref{2026山东02a} 所示，汽车在平直公路上行驶，路边有等间距的树木，车载摄像机记录了沿途景色。某同学根据一段视频绘制了图~\ref{2026山东02b}，横坐标 $N$ 为树木序号，纵坐标 $t$ 为对应树木出现的时刻，$t_{3} \sim t_{4}$ 内汽车通过隧道。关于汽车的运动，下列说法合理的是 \xzanswer{B}
\twopicture[numstyle=chinese, labelA=2026山东02a, labelB=2026山东02b, wA=7.5cm, wB=4.2cm]
{figs/02a.png}
{figs/02b.png}
\fourchoices[answer=B]
{$t_{1} \sim t_{2}$ 内做加速运动}
{$t_{2} \sim t_{3}$ 内做减速运动}
{$t_{1} \sim t_{2}$ 内的路程小于 $t_{2} \sim t_{3}$ 内的路程}
{$t_{2} \sim t_{3}$ 内的路程大于 $t_{4} \sim t_{5}$ 内的路程}
%% answer: B
%% memo:
\memoanswer{
树木是等间距 $d$ 的，故 $t-N$ 图像可类比 $t-x$ 图像。

AB．根据 $x = vt$ 可知，$t-x$ 图像斜率的倒数表示速度，由图知 $t_{1} \sim t_{3}$ 时间内斜率增大，则速度减小，做减速运动，故 A 错误，B 正确；

C．由图可知 $t_{1} \sim t_{2}$ 内的路程为 $4d$，大于 $t_{2} \sim t_{3}$ 内的路程 $3d$，故 C 错误；

D．$t_{2} \sim t_{3}$ 内的路程 $3d$，等于 $t_{4} \sim t_{5}$ 内的路程 $3d$，故 D 错误。

故选 B。
}

\item
%% number: 3
%% paperName: 2026年山东高考第 3 题 3 分
%% typeId: 1
%% type: 单选题
%% chapter: 气体性质
%% point: 分子动理论
%% method: 无
%% score: 3
%% degree: 650
%% duplicateId: 0
%% body:
一定质量的理想气体由状态 Ⅰ 变化到状态 Ⅱ，两种状态下气体分子的速率分布如图所示，图中 $f(v)$ 是速率 $v$ 附近单位速率区间内分子数占总分子数的百分比。下列说法正确的是 \xzanswer{D}
\onepicture[w=6cm]{\includesvg{figs/03}}
\fourchoices[answer=D]
{气体一定从外界吸收热量}
{气体中每个分子的速率都增加}
{速率 $v_{1}$ 附近单位速率区间内的分子数增加}
{气体中速率在 $v_{1} \sim v_{2}$ 区间的分子数占总分子数的比例减小}
%% answer: D
%% memo:
\memoanswer{
根据理想气体分子速率分布规律：温度越高，分子平均速率越大，速率分布曲线的峰值向速率更大的方向移动，且峰值降低、曲线更平缓。由此可知，图中状态 Ⅱ 温度高于状态 Ⅰ，气体从 Ⅰ 到 Ⅱ 温度升高。

A．一定质量理想气体的内能仅与温度有关，温度升高则内能增大 $\Delta U > 0$。根据热力学第一定律 $\Delta U = Q + W$，若外界对气体做功，则可能 $Q < 0$，气体可能是放热，A 错误；

B．温度升高是分子平均速率增大，属于统计规律，不代表每个分子的速率都增加，B 错误；

C．$f(v)$ 的物理意义是速率 $v$ 附近单位速率区间内分子数占总分子数的百分比，由图可知状态 Ⅱ 在 $v_{1}$ 处的 $f(v)$ 小于状态 Ⅰ，因此 $v_{1}$ 附近单位速率区间内的分子数减少，C 错误；

D．状态 Ⅱ 的峰值右移，速率小于 $v_{1}$ 的占比减少、速率大于 $v_{2}$ 的占比增加，由图可知，$v_{1} \sim v_{2}$ 区间图像面积减小，则气体中速率在 $v_{1} \sim v_{2}$ 区间的分子数占总分子数的比例减小，D 正确。

故选 D。
}

\item
%% number: 4
%% paperName: 2026年山东高考第 4 题 3 分
%% typeId: 1
%% type: 单选题
%% chapter: 曲线运动
%% point: 开普勒定律
%% method: 无
%% score: 3
%% degree: 650
%% duplicateId: 0
%% body:
海王星的卫星海卫二绕海王星的公转周期与地球公转周期近似相等。若太阳与海王星的质量比为 $a$，定义地球与太阳间的距离为 $1$ 个天文单位（$1\UAU$），则海卫二公转轨道的半长轴约为 \xzanswer{A}
\fourchoices[answer=A]
{$\frac{1}{\sqrt[3]{a}}\UAU$}
{$\frac{1}{\sqrt{a}}\UAU$}
{$\sqrt[3]{a}\UAU$}
{$\sqrt{a}\UAU$}
%% answer: A
%% memo:
\memoanswer{
地球绕太阳公转：设地球公转周期为 $T_{\text{地}}$，日地距离 $r_{\text{地}} = 1\UAU$，太阳质量为 $M_{\text{太}}$，则

\[ \frac{GM_{\text{太}}m_{\text{地}}}{r_{\text{地}}^{2}} = m_{\text{地}}\left(\frac{2\pi}{T_{\text{地}}}\right)^{2}r_{\text{地}} \]

海卫二绕海王星公转：设海卫二公转周期为 $T_{\text{卫}}$，轨道半长轴为 $r_{\text{卫}}$，海王星质量为 $M_{\text{海}}$，则

\[ \frac{GM_{\text{海}}m_{\text{卫}}}{r_{\text{卫}}^{2}} = m_{\text{卫}}\left(\frac{2\pi}{T_{\text{卫}}}\right)^{2}r_{\text{卫}} \]

已知 $T_{\text{卫}} \approx T_{\text{地}}$，且 $\frac{M_{\text{太}}}{M_{\text{海}}} = a$，则 $\frac{r_{\text{卫}}^{3}}{r_{\text{地}}^{3}} = \frac{M_{\text{海}}}{M_{\text{太}}} = \frac{1}{a}$

整理得 $r_{\text{卫}}^{3} = \frac{1}{a} \cdot (1\UAU)^{3}$

故 $r_{\text{卫}} = \frac{1}{\sqrt[3]{a}}\UAU$

故选 A。
}

\item
%% number: 5
%% paperName: 2026年山东高考第 5 题 3 分
%% typeId: 1
%% type: 单选题
%% chapter: 光学
%% point: 光的折射
%% method: 无
%% score: 3
%% degree: 550
%% duplicateId: 0
%% body:
半径为 $2\sqrt{2}\Ucm$ 的半圆形玻璃砖横截面如图所示，底部等腰梯形区域充满不反射、不透光的吸光性材料。梯形下底中点与圆心 $O$ 重合，下底为 $2\sqrt{2}\Ucm$，上底为 $2\Ucm$，高为 $1\Ucm$。单色光线从左侧入射，方向始终平行于直径 $AB$。当入射点从 $A$ 点逐渐向上移动至 $C$ 点时，光线恰能从玻璃砖另一侧射出，$C$ 点与 $AB$ 的距离为 $2\Ucm$，该玻璃的折射率为 \xzanswer{C}
\onepicture[w=6cm]{\includesvg{figs/05}}
\fourchoices[answer=C]
{$\sqrt{2}$}
{$\sqrt{5}$}
{$\frac{\sqrt{10}}{2}$}
{$\frac{\sqrt{5}}{2}$}
%% answer: C
%% memo:
\memoanswer{
根据题意作出光线进入玻璃砖过程的法线，标出入射角 $i$ 和折射角 $r$，如图所示

\onepicture[w=7cm]{figs/05_an.jpg}

由几何关系有

\[ \sin i = \frac{h_{1}}{R} = \frac{2\Ucm}{2\sqrt{2}\Ucm} = \frac{\sqrt{2}}{2} \]

可得 $i = 45^{\circ}$

则 $\alpha = 180^{\circ} - 2i = 90^{\circ}$

由几何关系可得 $OC = \frac{h_{2}}{\sin i} = \sqrt{2}\Ucm$

则 $\sin r = \frac{OC}{\sqrt{OC^{2} + R^{2}}} = \frac{\sqrt{2}}{\sqrt{10}} = \frac{\sqrt{5}}{5}$

则该玻璃的折射率为 $n = \frac{\sin i}{\sin r} = \frac{\sqrt{10}}{2}$

故选 C。
}

\item
%% number: 6
%% paperName: 2026年山东高考第 6 题 3 分
%% typeId: 1
%% type: 单选题
%% chapter: 功和能
%% point: 功率
%% method: 无
%% score: 3
%% degree: 650
%% duplicateId: 0
%% body:
“外骨骼机器人”是一种能增强运动能力的可穿戴装置。如图所示，在倾角为 $\theta$ 的斜坡上，某同学最多能拉着质量为 $m$ 的物体以恒定速度 $v$ 沿斜坡向上运动；穿戴“外骨骼机器人”后，最多能拉着质量为 $1.4m$ 的物体仍以相同的速度沿斜坡向上运动，绳子始终平行于斜坡，物体与斜坡之间的动摩擦因数为 $\mu$，重力加速度大小为 $g$，则该同学穿戴装置后，拉力的功率增加了 \xzanswer{A}
\onepicture[w=6cm]{\includesvg{figs/06}}
\fourchoices[answer=A]
{$0.4mgv(\sin\theta + \mu\cos\theta)$}
{$1.4mgv(\cos\theta + \mu\sin\theta)$}
{$0.4mgv(\sin\theta - \mu\cos\theta)$}
{$1.4mgv(\cos\theta - \mu\sin\theta)$}
%% answer: A
%% memo:
\memoanswer{
该同学拉着物体沿斜坡向上匀速运动时，拉力的功率等于物体克服重力与摩擦力做功的功率，则穿戴装置前有 $P_{1} = (mg\sin\theta + \mu mg\cos\theta)v$

穿戴装置后有 $P_{2} = (1.4mg\sin\theta + 1.4\mu mg\cos\theta)v$

所以拉力的功率的增加量为 $\Delta P = P_{2} - P_{1} = 0.4mgv(\sin\theta + \mu\cos\theta)$

故选 A。
}

\item
%% number: 7
%% paperName: 2026年山东高考第 7 题 3 分
%% typeId: 1
%% type: 单选题
%% chapter: 交流电
%% point: 交流电
%% method: 无
%% score: 3
%% degree: 450
%% duplicateId: 0
%% body:
图~\ref{2026山东07a} 是某交流发电机的原理图，在 $Oxy$ 坐标系中，以 $O$ 为圆心的上、下两个半圆区域内分别充满垂直 $Oxy$ 面向外和向里的匀强磁场，磁感应强度大小均为 $B_{1} = B_{0}$，金属棒一端始终位于 $O$ 点，绕 $O$ 点在 $Oxy$ 面内以恒定角速度 $\omega$ 转动，产生有效值为 $E_{\text{甲}}$、周期为 $T_{\text{甲}}$ 的交变电动势。如图~\ref{2026山东07b} 所示，现将一、四象限磁场方向分别变为垂直 $Oxy$ 面向里和向外，整个圆形区域磁感应强度大小变为 $B_{2} = 4B_{0}$，且金属棒转动角速度变为 $2\omega$，产生有效值为 $E_{\text{乙}}$、周期为 $T_{\text{乙}}$ 的交变电动势。则 \xzanswer{D}
\twopicture[numstyle=chinese, labelA=2026山东07a, labelB=2026山东07b, w=5cm]
{figs/07a.png}
{figs/07b.png}
\fourchoices[answer=D]
{$E_{\text{乙}} = 4E_{\text{甲}}$、$T_{\text{乙}} = \frac{1}{2}T_{\text{甲}}$}
{$E_{\text{乙}} = 8E_{\text{甲}}$、$T_{\text{乙}} = \frac{1}{2}T_{\text{甲}}$}
{$E_{\text{乙}} = 4E_{\text{甲}}$、$T_{\text{乙}} = \frac{1}{4}T_{\text{甲}}$}
{$E_{\text{乙}} = 8E_{\text{甲}}$、$T_{\text{乙}} = \frac{1}{4}T_{\text{甲}}$}
%% answer: D
%% memo:
\memoanswer{
根据题意，由图可知，图甲中金属棒转动 $2\pi$ 为一个周期，则有 $T_{\text{甲}} = \frac{2\pi}{\omega}$

图乙中金属棒转动 $\pi$ 为一个周期，则有 $T_{\text{乙}} = \frac{\pi}{2\omega}$

则有 $T_{\text{乙}} = \frac{1}{4}T_{\text{甲}}$

根据题意可知，甲、乙图中产生的交流电均为矩形交流电且正负最大电动势相等，则有

\[ E_{\text{甲}m} = B_{0}L \cdot \frac{0 + \omega L}{2} = \frac{1}{2}B_{0}\omega L^{2} \]

\[ E_{\text{乙}m} = 4B_{0}L \cdot \frac{0 + 2\omega L}{2} = 4B_{0}\omega L^{2} \]

作出图像，如图所示

\onepicture[w=7cm]{figs/07_an.jpg}

可知甲、乙图中产生的交流电的有效值为 $E_{\text{甲}} = E_{\text{甲}m} = \frac{1}{2}B_{0}\omega L^{2}$，$E_{\text{乙}} = E_{\text{乙}m} = 4B_{0}\omega L^{2}$

则有 $E_{\text{乙}} = 8E_{\text{甲}}$

故选 D。
}

\item
%% number: 8
%% paperName: 2026年山东高考第 8 题 3 分
%% typeId: 1
%% type: 单选题
%% chapter: 力物体的平衡
%% point: 三力平衡
%% method: 无
%% score: 3
%% degree: 450
%% duplicateId: 0
%% body:
如图所示，由光滑刚性杆组成的正四面体框架放置在水平面上，三条棱上各套有一个质量为 $m$ 的小球。三个小球通过相同的轻质弹簧连接，静止时恰好处于同一水平面。已知弹簧始终在弹性限度内，劲度系数为 $k$，重力加速度大小为 $g$，则每根弹簧的伸长量为 \xzanswer{B}
\onepicture[w=5cm]{figs/08.png}
\fourchoices[answer=B]
{$\frac{2\sqrt{3}mg}{3k}$}
{$\frac{6\sqrt{3}mg}{3k}$}
{$\frac{\sqrt{2}mg}{k}$}
{$\frac{\sqrt{3}mg}{k}$}
%% answer: B
%% memo:
\memoanswer{
根据题意可知，每根弹簧的伸长量相等，设每根弹簧的伸长量为 $\Delta x$，光滑刚性杆与地面的夹角为 $\theta$，结合正四面体的特点，由几何关系可得

\[ \sin\theta = \frac{\sqrt{6}}{3},\quad \cos\theta = \frac{\sqrt{3}}{3} \]

设两根弹簧对小球的弹力为 $F$，则有 $F = 2k\Delta x\cos 30^{\circ}$

对其中一个小球受力分析，如图所示

\onepicture[w=5cm]{figs/08_an.jpg}

沿杆方向上，由平衡条件有 $mg\sin\theta = F\cos\theta$

联立解得 $\Delta x = \frac{6\sqrt{3}mg}{3k}$

故选 B。
}

\item
%% number: 9
%% paperName: 2026年山东高考第 9 题 4 分
%% typeId: 2
%% type: 多选题
%% chapter: 机械波
%% point: 波的图象
%% method: 无
%% score: 4
%% degree: 450
%% duplicateId: 0
%% body:
波源 $O$ 在 $t = 0$ 时开始振动，产生一列沿 $x$ 轴正方向传播的简谐横波。$t = 4\Us$ 时质点 $M$ 刚好开始振动，质点 $M$ 与 $N$ 的平衡位置相距 $5\Um$。某时刻的波形如图所示，下列说法正确的是 \xzanswer{BD}
\onepicture[w=6cm]{\includesvg{figs/09}}
\fourchoices[answer=BD]
{波速为 $2\Ums$}
{波长为 $4\Um$}
{从图示状态经过 $2\Us$，$M$ 和 $N$ 速度相同}
{从图示状态经过 $3\Us$，$M$ 和 $N$ 偏离平衡位置的位移相同}
%% answer: BD
%% memo:
\memoanswer{
B．根据题图有 $\frac{5}{4}\lambda = 5\Um$

可得波长 $\lambda = 4\Um$，故 B 正确；

A．结合 B 项分析可知 $x_{OM} = \frac{1}{2}\lambda = 2\Um$

所以该波的波速 $v = \frac{x_{OM}}{t} = 0.5\Ums$，故 A 错误；

C．从题图示状态经过 $2\Us$，则波沿 $x$ 轴正方向传播了 $\Delta x = 1\Um = \frac{\lambda}{4}$，故 $M$ 从平衡位置运动到了波峰位置，速度为 $0$，$N$ 从波谷位置运动到了平衡位置，速度最大，显然两点速度不同，故 C 错误；

D．该波的周期 $T = \frac{\lambda}{v} = 8\Us$

从图示时刻计时，$M$ 的振动方程为 $y_{M} = A\sin\frac{\pi}{4}t$

$N$ 的振动方程为 $y_{N} = A\sin\left(\frac{\pi}{4}t - \frac{\pi}{2}\right)$

经 $t = 3\Us$ 后，$y_{M} = y_{N} = \frac{\sqrt{2}}{2}A$，故 D 正确。

故选 BD。
}

\item
%% number: 10
%% paperName: 2026年山东高考第 10 题 4 分
%% typeId: 2
%% type: 多选题
%% chapter: 电场
%% point: 带电粒子的圆周运动
%% method: 无
%% score: 4
%% degree: 450
%% duplicateId: 0
%% body:
融合新型功能材料的传感器在智能感知领域得到广泛应用。如图所示，光滑、绝缘椭圆轨道竖直放置，长轴 $AC = 2a$、短轴 $BD = 2b$，$AC$ 与 $BD$ 交于 $O$ 点，$B$ 为最低点，$A$ 点处内置可感知轨道压力的传感器。空间内充满平行轨道平面、斜向上的匀强电场（图中未画出）。质量为 $m$、电荷量为 $+q$ 的小球置于 $A$ 点时恰好静止，此时传感器示数等于 $mg$，$g$ 为重力加速度的大小。下列说法正确的是 \xzanswer{BCD}
\onepicture[w=8cm]{\includesvg{figs/10}}
\fourchoices[answer=BCD]
{电场强度的大小为 $\frac{mg}{q}$}
{电场强度方向与 $OA$ 的夹角为 $45^{\circ}$}
{A、B 两点间的电势差 $U_{AB} = \frac{(a + b)mg}{q}$}
{若小球在 B 点的速度水平向右、大小为 $v_{0}$，则到达 D 点时速度大小仍为 $v_{0}$}
%% answer: BCD
%% memo:
\memoanswer{
AB．作出小球在 A 点时的受力分析图如图 1 所示

\onepicture[w=4cm]{figs/10_an1.jpg}

由力的平衡条件结合几何关系有

\[ qE = \sqrt{(mg)^{2} + (mg)^{2}} = \sqrt{2}mg \]

\[ E = \frac{\sqrt{2}mg}{q} \]

所以电场强度的大小为 $E = \frac{\sqrt{2}mg}{q}$，且其方向与 $OA$ 的夹角 $\alpha$ 满足

$\tan\alpha = 1$，则 $\alpha = 45^{\circ}$，故 A 错误，B 正确；

C．分别作出 A、B 两点在经过 O 点的电场线上的两等势点 $A'$、$B'$，如图 2 所示

\onepicture[w=5cm]{figs/10_an2.jpg}

则由几何关系可知，A、B 两点间沿电场方向的距离为

\[ d = A'O + OB' = a\cos 45^{\circ} + b\cos 45^{\circ} = (a + b)\cos 45^{\circ} \]

所以 A、B 两点间的电势差为 $U_{AB} = -Ed = \frac{(a + b)mg}{q}$，故 C 正确；

D．小球所受支持力对其始终不做功，根据题意可知，小球所受电场力与重力的合力水平向右，$BD$ 连线沿竖直方向，所以从 B 到 D 的过程，小球所受合外力做的总功为 $0$，由动能定理可知其动能变化量为 $0$，故小球在 B、D 两点的速度大小均为 $v_{0}$，故 D 正确。

故选 BCD。
}

\item
%% number: 11
%% paperName: 2026年山东高考第 11 题 4 分
%% typeId: 2
%% type: 多选题
%% chapter: 运动和力
%% point: 动量和能量
%% method: 无
%% score: 4
%% degree: 350
%% duplicateId: 0
%% body:
如图所示，质量相等的两个小物块 M 和 N，M 恰好静止于倾角为 $\theta$ 的固定斜面上，N 从斜面上某位置由静止释放，$t_{0}$ 时刻以速度 $v$ 与 M 发生弹性碰撞。已知 M 与斜面间动摩擦因数为 $\tan\theta$，最大静摩擦力等于滑动摩擦力，N 与斜面间无摩擦，碰撞时间极短，斜面足够长，下列描述 M、N 速度规律的 $v_{M}-t$、$v_{N}-t$ 图像正确的是 \xzanswer{AD}
\onepicture[w=6cm]{\includesvg{figs/11a}}
\fourchoices[answer=AD, ispicture=true, h=2.7cm]
{figs/11b.png}
{figs/11c.png}
{figs/11d.png}
{figs/11e.png}
%% answer: AD
%% memo:
\memoanswer{
M、N 质量相等，根据动量守恒定律和能量守恒定律有

\[ mv = mv_{M} + mv_{N},\quad \frac{1}{2}mv^{2} = \frac{1}{2}mv_{M}^{2} + \frac{1}{2}mv_{N}^{2} \]

可得碰撞后 $v_{M} = v$，$v_{N} = 0$

即发生弹性碰撞后交换速度，则 $t_{0}$ 时刻两物块碰后瞬间 M 的速度为 $v$，N 的速度为 $0$，M 与斜面间的动摩擦因数为 $\tan\theta$、N 与斜面间无摩擦，则碰后 M 做匀速直线运动，N 做加速度大小为 $a = \frac{v}{t_{0}}$ 的匀加速直线运动，从 M、N 第一次碰撞后瞬间到第二次碰撞前瞬间的过程，两物块位移相等，设该过程运动时间为 $\Delta t_{1}$，则有

\[ v\Delta t_{1} = \frac{1}{2}a(\Delta t_{1})^{2} \]

解得 $\Delta t_{1} = 2t_{0}$

则 $3t_{0}$ 时刻 M、N 发生第二次碰撞，碰前瞬间 M 的速度为 $v$，N 的速度为 $v' = a\Delta t_{1} = 2v$

M、N 质量相等，发生弹性碰撞后交换速度，所以第二次碰后瞬间 M 的速度为 $v' = 2v$，N 的速度为 $v$，同理可得 M、N 从第二次碰撞后瞬间到第三次碰撞前瞬间有

\[ v'\Delta t_{2} = v\Delta t_{2} + \frac{1}{2}a(\Delta t_{2})^{2} \]

解得 $\Delta t_{2} = 2t_{0}$

所以 $5t_{0}$ 时刻发生第三次碰撞，第三次碰前瞬间 M 的速度为 $v' = 2v$，N 的速度为 $v'' = v + a\Delta t_{2} = 3v$

对比选项中图像可知 AD 正确，BC 错误。

故选 AD。
}

\item
%% number: 12
%% paperName: 2026年山东高考第 12 题 4 分
%% typeId: 2
%% type: 多选题
%% chapter: 电磁感应
%% point: 电磁感应含容电路
%% method: 无
%% score: 4
%% degree: 350
%% duplicateId: 0
%% body:
如图所示，足够长 U 形光滑金属导轨固定在绝缘水平面上，宽度为 $L$，电阻不计。区域 Ⅰ 为正方形，充满垂直轨道平面向上、磁感应强度大小 $B_{1}$ 随时间 $t$ 均匀变化的匀强磁场，即 $B_{1} = kt$（$k$ 为大于零的常量）；区域 Ⅱ 内，导轨上接有开关 $\text{S}_{1}$、$\text{S}_{2}$，导轨间接有电容为 $C_{1}$、$C_{2}$ 的两电容器；区域 Ⅲ 内充满方向垂直轨道平面向下、磁感应强度大小为 $B$ 的匀强磁场。初始时，质量为 $m$、有一定阻值的导体棒 MN 静止于区域 Ⅲ 中某处，$\text{S}_{1}$ 闭合，$\text{S}_{2}$ 断开，$C_{1}$ 充电。$C_{1}$ 充电完毕后，断开 $\text{S}_{1}$，闭合 $\text{S}_{2}$，MN 开始运动，经过一段时间系统达到最终稳定状态。MN 长度与导轨宽度相同并始终与导轨接触良好。下列说法正确的是 \xzanswer{AC}
\onepicture[w=7cm]{\includesvg{figs/12}}
\fourchoices[answer=AC]
{$C_{1}$ 充电完毕时的电荷量为 $C_{1}kL^{2}$}
{$C_{1}$ 充电完毕时，上极板带负电}
{最终 MN 的速度为 $\frac{kC_{1}BL^{3}}{B^{2}L^{2}(C_{1} + C_{2}) + m}$}
{最终 $C_{2}$ 的电荷量为 $\frac{kC_{1}C_{2}B^{2}L^{4}}{B^{2}L^{2}C_{2} + m}$}
%% answer: AC
%% memo:
\memoanswer{
A．$\text{S}_{1}$ 闭合、$\text{S}_{2}$ 断开时，区域 Ⅰ 中产生感应电动势给 $C_{1}$ 充电，由法拉第电磁感应定律可知产生的感应电动势为 $U_{1} = \frac{\Delta B_{1}}{\Delta t}L^{2} = kL^{2}$

由电容定义式可得 $C_{1}$ 充电完毕时的电荷量为 $Q_{1} = C_{1}U_{1} = C_{1}kL^{2}$，故 A 正确；

B．由楞次定律和右手螺旋定则可知区域 Ⅰ 中产生顺时针方向的感应电流，则 $C_{1}$ 充电完毕时，上极板带正电，故 B 错误；

CD．断开 $\text{S}_{1}$、闭合 $\text{S}_{2}$ 时，$C_{1}$ 放电，最终回路中感应电流为 $0$，$C_{1}$、$C_{2}$ 和 MN 两端电压相等，设最终稳定时 $C_{1}$ 的电荷量为 $Q_{1}'$，$C_{2}$ 的电荷量为 $Q_{2}$，流过 MN 的电荷量为 $Q_{MN}$，MN 的速度为 $v$，对 MN 由动量定理和电流的定义式 $q = It$ 有 $B\bar{I}Lt = BLQ_{MN} = mv$

结合法拉第电磁感应定律和电容的定义式可知 $\frac{Q_{1}'}{C_{1}} = \frac{Q_{2}}{C_{2}} = BLv$

联立可得 $Q_{MN} = \frac{m}{BL} \cdot \frac{Q_{2}}{C_{2}BL} = \frac{mQ_{2}}{C_{2}B^{2}L^{2}}$，$Q_{1}' = \frac{C_{1}}{C_{2}}Q_{2}$

整个过程由电荷守恒定律有 $Q_{1} = Q_{1}' + Q_{2} + Q_{MN}$

联立可得 $Q_{2} = \frac{kC_{1}C_{2}B^{2}L^{4}}{B^{2}L^{2}(C_{1} + C_{2}) + m}$，$v = \frac{Q_{2}}{C_{2}BL} = \frac{kC_{1}BL^{3}}{B^{2}L^{2}(C_{1} + C_{2}) + m}$，故 C 正确，D 错误。

故选 AC。
}

\item
%% number: 13
%% paperName: 2026年山东高考第 13 题 6 分
%% typeId: 4
%% type: 实验题
%% chapter: 力物体的平衡
%% point: 摩擦力
%% method: 无
%% score: 6
%% degree: 650
%% duplicateId: 0
%% body:
某同学利用智能小车测量物体间动摩擦因数。如图~\ref{2026山东13a} 所示，将橡胶板平整固定在水平桌面上，把质量为 $200\Ug$ 的物块放置在橡胶板的右端，用细线与小车连接，调整细线与橡胶板平行。
\onepicture[num=甲, label=2026山东13a, align=right, w=8cm]{\includesvg{figs/13a}}

智能小车启动，缓慢拉动物块，使物块由静止逐渐变为匀速直线运动，小车每隔 $0.2\Us$ 采集一组拉力 $F$ 与时间 $t$ 的数据，实验数据如图~\ref{2026山东13b} 所示。回答以下问题：
\onepicture[num=乙, label=2026山东13b, align=right, w=6cm]{figs/13b.png}
\begin{enumerate}
	\item
	在 $t < 2.6\Us$ 范围内，物块处于 \tkanswer{静止}（填“静止”或“匀速直线运动”）状态。
	\item
	为得到物块运动过程中的滑动摩擦力，应选用时间段 \tkanswer{$4.0 \sim 6.0$} $\Us$ 内的数据。（填“$2.8 \sim 6.0$”或“$4.0 \sim 6.0$”）
	\item
	当地重力加速度大小为 $9.8\Umsq$，此物块与橡胶板间动摩擦因数为 \tkanswer{$0.56$}（保留两位有效数字）。
	\item
	若桌面不水平，右端略高于左端，则会导致动摩擦因数的测量值比实际值 \tkanswer{偏小}（填“偏大”或“偏小”）。
\end{enumerate}
%% answer:
\jdanswer{
\begin{enumerate}
	\item
	静止

	\item
	$4.0 \sim 6.0$

	\item
	$0.56$

	\item
	偏小

\end{enumerate}
}
%% memo:
\memoanswer{
（1）由于智能小车缓慢拉动物块使其由静止逐渐变为匀速直线运动，因此在初始一段时间内，物块处于静止状态。

（2）只有当物块匀速运动时，细线上拉力 $F$ 才与物块受到的滑动摩擦力大小相等，因此需要选用物块匀速运动阶段的数据，即题图乙中时间段 $4.0 \sim 6.0\Us$ 内的数据。

（3）根据题图乙可知，物块受到的滑动摩擦力 $f = F = 1.10\UN$

根据 $f = \mu F_{N}$，结合 $F_{N} = mg$，解得 $\mu = \frac{F}{mg} \approx 0.56$

（4）若桌面右端略高于左端，橡胶板形成一个左低右高的斜面。物块初始放在右端，被小车向左缓慢拉动，属于沿斜面向下运动。此时物块受到重力沿斜面向下的分力 $mg\sin\theta$ 与拉力 $F'$ 同向，共同克服滑动摩擦力 $\mu mg\cos\theta$，即 $F' + mg\sin\theta = \mu mg\cos\theta$

得 $F' = \mu mg\cos\theta - mg\sin\theta$

若仍按水平面的公式 $\mu_{\text{测}} = \frac{F'}{mg}$ 计算，则 $\mu_{\text{测}} = \mu\cos\theta - \sin\theta$

由于倾角 $\theta > 0$，$\cos\theta < 1$，$\sin\theta > 0$，故 $\mu_{\text{测}} < \mu$，即测量值偏小。
}

\item
%% number: 14
%% paperName: 2026年山东高考第 14 题 8 分
%% typeId: 4
%% type: 实验题
%% chapter: 光学
%% point: 双缝干涉测量波长
%% method: 无
%% score: 8
%% degree: 550
%% duplicateId: 0
%% body:
某实验小组用光传感器做双缝干涉实验并测量光的波长 $\lambda$。
\begin{enumerate}
	\item
	如图~\ref{2026山东14a} 所示，将激光光源、间距 $d = 0.50\Umm$ 的双缝和光传感器依次安装在光具座上，打开光源，调整光路，此过程中应避免激光直接射入眼睛，因为激光具有 \tkanswer{亮度高}（填“亮度高”或“相干性好”）的特点。
	\onepicture[num=甲, label=2026山东14a, align=right, w=8cm]{figs/14a.png}
	\item
	调整光路后，各器件位置如图~\ref{2026山东14a} 所示，双缝到光传感器的距离 $l =$ \tkanswer{$31.25$} $\Ucm$。
	\item
	连接光传感器至计算机，得到干涉条纹的光强分布如图~\ref{2026山东14b} 所示。测得 A 区域的宽度为 $2.40\Umm$，则条纹间距 $\Delta x$ \tkanswer{$0.40$} $\Umm$（保留两位有效数字）。
	\onepicture[num=乙, label=2026山东14b, align=right, w=6cm]{figs/14b.png}
	\item
	利用公式 $\lambda =$ \tkanswer{$\frac{d}{l}\Delta x$}（用 $l$、$d$ 和 $\Delta x$ 表示），某同学计算得该激光的波长 $\lambda$ 为 $640\Unm$，此波长的光为电磁波中的 \tkanswer{可见光}（填“红外线”“可见光”或“紫外线”）。
\end{enumerate}
%% answer:
\jdanswer{
\begin{enumerate}
	\item
	亮度高

	\item
	$31.25$

	\item
	$0.40$

	\item
	$\frac{d}{l}\Delta x$，可见光

\end{enumerate}
}
%% memo:
\memoanswer{
（1）激光的亮度高特点使它可以在很小的空间和很短的时间内集中很大的能量，如果直接射入眼睛会对眼睛造成不可逆转的伤害。

（2）根据题图甲可知，双缝所在位置为 $x_{1} = 12.00\Ucm$，光传感器所在的位置为 $x_{2} = 43.25\Ucm$，因此双缝到光传感器的距离 $l = x_{2} - x_{1} = 31.25\Ucm$

（3）根据题图乙可知，A 区域有 $6$ 个条纹间距，则条纹间距 $\Delta x = \frac{x_{A}}{6} = 0.40\Umm$

（4）根据双缝干涉条纹间距公式

\[ \Delta x = \frac{l}{d}\lambda \]

得 $\lambda = \frac{d}{l}\Delta x$

波长 $\lambda = 640\Unm$ 在可见光波段，因此此波长的光为电磁波中的可见光。
}

\item
%% number: 15
%% paperName: 2026年山东高考第 15 题 8 分
%% typeId: 6
%% type: 计算题
%% chapter: 气体性质
%% point: 理想气体状态方程
%% method: 无
%% score: 8
%% degree: 550
%% duplicateId: 0
%% body:
竖直固定的圆柱形透明管深度为 $l$，管内横截面积为 $S$；圆柱形物块长为 $\frac{1}{9}l$，横截面积为 $S$，密度为 $\rho$。室温 $T_{1} = 300\UK$ 时，某同学将表面涂润滑油的物块竖直置于管口封住管内气体，并使物块缓慢进入透明管，过程中气体无泄漏。当物块处于静止状态时，其上表面恰好与管口齐平，如图~\ref{2026山东15b} 所示。已知透明管与物块均具有良好导热性能，不计物块与透明管间的摩擦，重力加速度大小为 $g$，大气压强恒定，空气可视为理想气体。
\begin{enumerate}
	\item
	求当地大气压强 $p_{0}$；
	\item
	将装置放置较长时间后，物块下方气柱高度为 $\frac{7}{10}l$，该同学认为此装置漏气，测得此时室温 $T_{2} = 270\UK$，求管内剩余气体与密封刚完成时气体的质量比。
\end{enumerate}
\twopicture[numstyle=chinese, labelA=2026山东15a, labelB=2026山东15b, h=6cm, align=right]
{figs/15a.png}
{figs/15b.png}
%% answer:
\jdanswer{
\begin{enumerate}
	\item
	$p_{0} = \frac{8}{9}\rho gl$

	\item
	$\frac{7}{8}$

\end{enumerate}
}
%% memo:
\memoanswer{
（1）设密封刚完成时管内气体压强为 $p_{1}$，气柱长度为 $l - \frac{1}{9}l = \frac{8}{9}l$，体积 $V_{1} = \frac{8}{9}lS$

物块受力平衡有 $p_{1}S = p_{0}S + \rho \cdot \frac{1}{9}lS \cdot g$

解得 $p_{1} = p_{0} + \frac{1}{9}\rho gl$

初始时管内气体体积为 $V_{0} = lS$，压强为大气压 $p_{0}$，温度均为 $T_{1}$。过程等温，由玻意耳定律 $p_{0}lS = p_{1} \cdot \frac{8}{9}lS$

解得 $p_{0} = \frac{8}{9}\left(p_{0} + \frac{1}{9}\rho gl\right)$

化简有 $p_{0} = \frac{8}{9}\rho gl$

（2）装置放置较长时间后，物块下方气柱高度为 $\frac{7}{10}l$，体积 $V_{2} = \frac{7}{10}lS$，温度 $T_{2} = 270\UK$

物块受力平衡不变，故此时气体压强 $p_{2} = p_{0} + \frac{1}{9}\rho gl = \rho gl$

密封刚完成时状态 $p_{1} = \rho gl$，$V_{1} = \frac{8}{9}lS$，$T_{1} = 300\UK$

由理想气体状态方程，气体质量比等于 $\frac{pV}{T}$ 之比，且压强 $p$ 不变，则

\[ \frac{m_{2}}{m_{1}} = \frac{V_{2}}{V_{1}} \cdot \frac{T_{1}}{T_{2}} = \frac{\frac{7}{10}lS}{\frac{8}{9}lS} \cdot \frac{300}{270} = \frac{7}{8} \]
}

\item
%% number: 16
%% paperName: 2026年山东高考第 16 题 8 分
%% typeId: 6
%% type: 计算题
%% chapter: 电磁感应
%% point: 电磁感应受力分析
%% method: 无
%% score: 8
%% degree: 550
%% duplicateId: 0
%% body:
如图所示，平行金属导轨间距为 $L$，导轨平面与水平面夹角为定值，二者交线与导轨垂直。电动势均为 $E$、内阻为 $r$ 的两电源，开关 $\text{S}_{1}$、$\text{S}_{2}$ 及滑动变阻器 $R_{P}$ 与导轨相连，导轨处于磁感应强度为 $B$、方向平行于导轨向下的匀强磁场中。$\text{S}_{1}$ 闭合、$\text{S}_{2}$ 断开，质量为 $m$，长为 $L$，内阻为 $r$ 的金属杆 $cd$ 静止于导轨上。调节滑片，当 $R_{P}$ 阻值为 $r$ 时，$cd$ 刚好要下滑。保持滑片位置不变，断开 $\text{S}_{1}$，闭合 $\text{S}_{2}$，$cd$ 开始下滑并始终与导轨接触良好。已知 $cd$ 与导轨间动摩擦因数为 $\mu$，最大静摩擦力等于滑动摩擦力，不计导轨电阻。$cd$ 下滑过程中，求：
\begin{enumerate}
	\item
	$cd$ 的电功率 $P$；
	\item
	$cd$ 的加速度大小 $a$。
\end{enumerate}
\onepicture[w=7cm, align=right]{\includesvg{figs/16}}
%% answer:
\jdanswer{
\begin{enumerate}
	\item
	$P = \frac{E^{2}}{9r}$

	\item
	$a = \frac{2\mu BEL}{3mr}$

\end{enumerate}
}
%% memo:
\memoanswer{
（1）断开 $\text{S}_{1}$，闭合 $\text{S}_{2}$ 时回路电流 $I = \frac{E}{3r}$

则 $cd$ 的电功率 $P = I^{2}r$

解得 $P = \frac{E^{2}}{9r}$

（2）设斜面倾角为 $\theta$，$\text{S}_{1}$ 闭合、$\text{S}_{2}$ 断开 $cd$ 刚好要下滑时有 $mg\sin\theta = \mu(mg\cos\theta + BIL)$

由于两电源电动势均为 $E$、内阻为 $r$，则回路电流大小不变均为 $I = \frac{E}{3r}$

断开 $\text{S}_{1}$，闭合 $\text{S}_{2}$ 有 $mg\sin\theta - \mu(mg\cos\theta - BIL) = ma$

联立解得 $a = \frac{2\mu BEL}{3mr}$
}

\item
%% number: 17
%% paperName: 2026年山东高考第 17 题 14 分
%% typeId: 6
%% type: 计算题
%% chapter: 曲线运动
%% point: 圆周运动的多解
%% method: 无
%% score: 14
%% degree: 450
%% duplicateId: 0
%% body:
如图~\ref{2026山东17a} 所示，半径 $r = 1\Um$ 的圆筒竖直放置，上下底面圆心分别为 $O$ 和 $O'$，筒内有三个互成 $120^{\circ}$ 角且可以绕 $OO'$ 转动的竖直矩形叶片 $\text{S}_{1}$、$\text{S}_{2}$ 和 $\text{S}_{3}$，叶片与圆筒上下齐平，宽度等于圆筒半径，$N$ 为圆筒上底边沿一点。圆筒上底面固定一半径 $R = \frac{3}{\pi}\Um$ 的水平圆形轨道，轨道上有一长度可忽略的缺口位于 $O$ 点，$OM$ 为轨道直径且 $OM \perp ON$；在轨道的 $M$ 点放置一小物体，某时刻该物体在内力作用下突然分成 A、B 两部分并弹开，其质量分别为 $m_{A}$、$m_{B}$。A 以 $v_{A} = 1\Ums$ 的速率沿轨道运动至缺口进入圆筒，在筒内做平抛运动后恰好紧贴圆筒下底面沿边穿出。已知 $g = 10\Umsq$。忽略空气阻力及一切摩擦，轨道、圆筒和叶片的厚度均忽略不计。
\begin{enumerate}
	\item
	求 A 在轨道内运动的时间及圆筒高度；
	\item
	若叶片以恒定角速度 $\omega$ 顺时针转动，且 A 运动至缺口时，$\text{S}_{1}$ 恰好转过 $ON$ 位置，如图~\ref{2026山东17b} 所示，随后 A 未与叶片碰撞，从圆筒下底面沿边穿出时，$\text{S}_{2}$ 恰好转至 $ON$ 位置。
	\begin{enumerate}
		\item
		求角速度的大小 $\omega$；
		\item
		若 B 运动至缺口后能从任意两个叶片间的区域穿过圆筒，且未与叶片及筒壁碰撞，求 $\frac{m_{A}}{m_{B}}$。
	\end{enumerate}
\end{enumerate}
\twopicture[numstyle=chinese, labelA=2026山东17a, labelB=2026山东17b, h=6.5cm, align=right]
{figs/17a.png}
{figs/17b.png}
%% answer:
\jdanswer{
\begin{enumerate}
	\item
	$t_{0} = 3\Us$，$h = 5\Um$

	\item
	（i）$\omega = \frac{2\pi}{3}\Urads$，（ii）$\frac{m_{A}}{m_{B}} = \frac{6}{7 + 2k}$（$k = 0, 1, 2, 3, \ldots$）

\end{enumerate}
}
%% memo:
\memoanswer{
（1）A 在轨道内做匀速圆周运动，运动时间 $t_{0} = \frac{s}{v_{A}} = 3\Us$

由题知，A 以 $v_{A} = 1\Ums$ 的速率沿轨道运动至缺口进入圆筒，在筒内做平抛运动后恰好紧贴圆筒下底面沿边穿出，则在水平方向有 $r = v_{A}t$

解得 $t = 1\Us$

竖直方向有 $h = \frac{1}{2}gt^{2} = 5\Um$

（2）（i）由于 A 未与叶片碰撞，从圆筒下底面沿边穿出时，叶片转动了 $\theta = \frac{2\pi}{3}$，则叶片角速度 $\omega = \frac{\theta}{t} = \frac{2\pi}{3}\Urads$

（ii）物体在内力作用下突然分成 A、B 两部分并弹开有 $m_{A}v_{A} = m_{B}v_{B}$

则 $\frac{m_{A}}{m_{B}} = v_{B}$

由于 B 未与筒壁碰撞则有 $x_{B} = v_{B} \cdot 1 = \frac{m_{A}}{m_{B}} \leq r = 1\Um$

且 A 运动至缺口时，$\text{S}_{1}$ 恰好转过 $ON$ 位置，则以此时为时间基准，则 B 从 M 到 O 所用的时间

\[ \Delta t = \frac{\pi R}{v_{B}} - \frac{\pi R}{v_{A}} = 3\left(\frac{1}{v_{B}} - 1\right) \]

要使 B 运动至缺口后能从任意两个叶片间的区域穿过圆筒，则 B 运动到 O 时 $\text{S}_{1}$ 或 $\text{S}_{2}$ 或 $\text{S}_{3}$ 要运动到 $ON$ 反向的位置，则叶片转过的角度 $\theta' = \frac{\pi}{3} + \frac{2\pi}{3} \cdot k$（$k = 0, 1, 2, 3, \ldots$）

B 不与叶片碰撞必须满足 $\Delta t = \frac{\theta'}{\omega}$

解得 $v_{B} = \frac{6}{7 + 2k}\Ums$（$k = 0, 1, 2, 3, \ldots$）

则 $\frac{m_{A}}{m_{B}} = \frac{6}{7 + 2k}$（$k = 0, 1, 2, 3, \ldots$）
}

\item
%% number: 18
%% paperName: 2026年山东高考第 18 题 16 分
%% typeId: 6
%% type: 计算题
%% chapter: 磁场
%% point: 带电粒子在磁场中的运动
%% method: 无
%% score: 16
%% degree: 350
%% duplicateId: 0
%% body:
在 $xOy$ 坐标系中，第二象限有一粒子发生器，其右侧放置速度选择器，速度选择器中电场强度大小为 $E$，方向沿 $y$ 轴负方向，匀强磁场方向垂直 $xOy$ 面向里；$y = -x$（$x \leq 0$）与 $y$ 轴正半轴所围区域 Ⅰ 中充满垂直 $xOy$ 面向外的匀强磁场；$x$ 轴下方为区域 Ⅱ、第一象限为区域 Ⅲ，两区域均充满方向垂直 $xOy$ 面向里的匀强磁场，磁感应强度大小分别为 $5B$ 和 $12B$。质量为 $m$，电荷量为 $q$ 的粒子 a 经速度选择器后以速率 $v$ 从点（$-h$，$h$）沿 $x$ 轴正方向进入区域 Ⅰ，一段时间后恰好从原点 $O$ 沿 $y$ 轴负方向进入区域 Ⅱ。不计粒子重力及粒子间相互作用。
\begin{enumerate}
	\item
	求速度选择器中磁感应强度大小 $B_{0}$ 和区域 Ⅰ 中磁感应强度大小 $B_{1}$；
	\item
	求粒子 a 从 $O$ 点运动到 $P$ 点（$\frac{19mv}{30qB}$，0）的时间 $t$；
	\item
	当粒子 a 从点（$\frac{2mv}{5qB}$，0）离开区域 Ⅱ 进入区域 Ⅲ 时，和 a 电荷量相同的粒子 b 恰好从 $O$ 点以速率 $v$ 沿 $y$ 轴负方向进入区域 Ⅱ，若粒子 a、b 在 $x$ 轴相遇且相遇时速度都沿 $y$ 轴正方向，求粒子 b 的质量 $M$ 及第一次相遇时的 $x$ 轴坐标 $x_{1}$。
\end{enumerate}
\onepicture[w=8cm, align=right]{\includesvg{figs/18}}
%% answer:
\jdanswer{
\begin{enumerate}
	\item
	$B_{0} = \frac{E}{v}$，$B_{1} = \frac{mv}{qh}$

	\item
	$\frac{29\pi m}{60qB}$

	\item
	$M = \frac{17}{10}m$，$x_{1} = \frac{34mv}{15qB}$

\end{enumerate}
}
%% memo:
\memoanswer{
（1）粒子在速度选择器中做直线运动，有 $qE = qvB_{0}$

解得 $B_{0} = \frac{E}{v}$

粒子在区域 Ⅰ 中做匀速圆周运动，由几何关系得轨迹半径 $R_{1} = h$，根据 $qvB_{1} = m\frac{v^{2}}{R_{1}}$

解得 $B_{1} = \frac{mv}{qh}$

（2）根据 $qvB = m\frac{v^{2}}{R}$ 可知，粒子在区域 Ⅱ、Ⅲ 中的轨迹半径分别为

\[ R_{2} = \frac{mv}{5qB},\quad R_{3} = \frac{mv}{12qB} \]

画出粒子运动轨迹图

\onepicture[w=9cm]{figs/18_an1.jpg}

粒子 a 轨迹：$O \to A \to B \to P$

可知从 $O$ 点进入区域 Ⅱ，经半圆到达 $A\left(\frac{2mv}{5qB}, 0\right)$

再根据 $qvB = m\frac{4\pi^{2}}{T^{2}}R$，可知 $O$ 到 $A$ 运动时间 $t_{1} = \frac{\pi m}{5qB}$

从 $A$ 点进入区域 Ⅲ，经半圆到达 $B\left(\frac{7mv}{30qB}, 0\right)$，时间 $t_{2} = \frac{\pi m}{12qB}$

从 $B$ 点再次进入区域 Ⅱ，经半圆到达 $P\left(\frac{19mv}{30qB}, 0\right)$，时间 $t_{3} = \frac{\pi m}{5qB}$

则总时间 $t = t_{1} + t_{2} + t_{3} = \frac{29\pi m}{60qB}$

（3）画出粒子 a、b 的运动轨迹如图

\onepicture[w=11cm]{figs/18_an2.jpg}

区域 Ⅱ：a 粒子半径 $r_{2} = \frac{mv}{5qB}$，b 粒子半径 $\frac{Mv}{5qB}$，a 粒子半周期 $\tau_{2} = \frac{\pi m}{5qB}$，b 粒子半周期 $\frac{\pi M}{5qB}$

区域 Ⅲ：a 粒子半径 $r_{3} = \frac{mv}{12qB}$，b 粒子半径 $\frac{Mv}{12qB}$，a 粒子半周期 $\tau_{3} = \frac{\pi m}{12qB}$，b 粒子半周期 $\frac{\pi M}{12qB}$

粒子每完成一次“区域 Ⅱ 下半圆 ＋ 区域 Ⅲ 上半圆”的运动，$x$ 坐标的净变化 $\Delta x_{a} = 2r_{2} - 2r_{3} = \frac{7mv}{30qB}$

同理 $\Delta x_{b} = \frac{7Mv}{30qB}$

所花时间 $\Delta t_{a} = \tau_{2} + \tau_{3} = \frac{17\pi m}{60qB}$

同理 $\Delta t_{b} = \frac{17\pi M}{60qB}$

取粒子 a 从 $O$ 点沿 $y$ 轴负方向进入区域 Ⅱ 的时刻为 $t = 0$，粒子 a 第 1 次速度向上点：在区域 Ⅱ 中走下半圆到达 $x$ 轴，坐标

\[ x_{a,1} = 2r_{2} = \frac{2}{5} \cdot \frac{mv}{qB} = \frac{12mv}{30qB} \]

时间

\[ t_{a,1} = \tau_{2} = \frac{\pi m}{5qB} = \frac{12\pi m}{60qB} \]

此后每经过一次区域 Ⅲ 上半圆和一次区域 Ⅱ 下半圆，到达下一个速度向上点。设第 $n$ 次速度向上点 $n = 1, 2, 3, \ldots$，则

\[ x_{a,n} = x_{a,1} + (n - 1)\Delta x_{a} = \frac{12mv}{30qB} + (n - 1)\frac{7mv}{30qB} = \frac{7n + 5}{30} \cdot \frac{mv}{qB} \]

\[ t_{a,n} = t_{a,1} + (n - 1)\Delta t_{a} = \frac{12\pi m}{60qB} + (n - 1)\frac{17\pi m}{60qB} = \frac{17n - 5}{60} \cdot \frac{\pi m}{qB} \]

粒子 b 在 a 到达第一次向上点 a（即 $x = \frac{2mv}{5qB}$）时，从 $O$ 点以同样速率沿 $y$ 轴负方向进入区域 Ⅱ。b 的起始时刻为

\[ t_{a,1} = \frac{\pi m}{5qB} = \frac{12\pi m}{60qB} \]

粒子 b 第 1 次速度向上点（绝对时刻）：b 先在区域 Ⅱ 走下半圆，坐标

\[ x_{b,1} = 2 \cdot \frac{Mv}{5qB} = \frac{12Mv}{30qB} \]

时刻

\[ T_{b,1} = t_{a,1} + \frac{\pi M}{5qB} = \frac{12\pi m}{60qB} + \frac{12\pi M}{60qB} = \frac{12m + 12M}{60qB}\pi \]

同理，第 $k$ 次速度向上点 $k = 1, 2, 3, \ldots$

\[ x_{b,k} = \frac{12Mv}{30qB} + (k - 1)\frac{7Mv}{30qB} = \frac{7k + 5}{30} \cdot \frac{Mv}{qB} \]

\[ T_{b,k} = T_{b,1} + (k - 1)\left(\frac{\pi M}{5qB} + \frac{\pi M}{12qB}\right) = \frac{12m + (17k - 5)M}{60qB}\pi \]

由相遇条件 $x_{a,n} = x_{b,k}$、$t_{a,n} = t_{b,k}$ 得

\[ (7n + 5)m = (7k + 5)M \]

\[ (17n - 5)m = 12m + (17k - 5)M \]

两式相除消去 $M$ 得 $10n - 17k = 5$，最小正整数解为 $k = 5$，$n = 9$

代入得 $M = \frac{17}{10}m$，$x_{1} = \frac{34mv}{15qB}$
}

\end{enumerate}

\end{document}
